\documentclass[10pt]{mypackage}

\usepackage[uprightgreeks,noDcommand,nofligatures]{kpfonts}
\renewcommand{\coloneq}{\coloneqq}
\renewcommand{\eqcolon}{\eqqcolon}
%\usepackage{dsfont}
%\renewcommand{\mathbb}{\mathds}
%\usepackage{homework}
\usepackage{notes}

\usepackage[ backend=bibtex, style = alphabetic, sorting=ynt ]{biblatex}
\addbibresource{../all_references.bib}

\usepackage{parskip}
\frenchspacing

\fancyhf{}
\fancyhead[R]{Avinash Iyer}
\fancyhead[L]{Factor Maps and the Point Map Realization Theorem}
\fancyfoot[C]{\thepage}

\setcounter{secnumdepth}{0}

\begin{document}
\RaggedRight
\section{Polish Spaces and Borel Spaces}%
\begin{definition}
  A space topological space $\left( X,\mathcal{T} \right)$ is called a Polish space if $X$ is separable and there is a compatible complete metric.
\end{definition}
\begin{theorem}
  Let $X$ be an uncountable Polish space. Then, $X$ contains a Cantor set.
\end{theorem}
\begin{proof}
  Let $d$ be a compatible complete metric. Let $Z$ be the set of all points $x\in X$ such that each open neighborhood of $x$ is uncountable. Then, $Z$ is closed.

  For each $x\in Z$, let $r_x$ be the supremum of all $0 < r < 1$ such that the open ball centered at $x$ with radius $r$ is countable. Let $S$ be a countable dense subset of $X\setminus Z$. The open balls around $x\in S$ with radius $\frac{1}{2}r_x$ for each $x\in S$ cover $X\setminus Z$. S, $X\setminus Z$ is countable.

  Assuming $Z = X$, then by a standard recursive procedure, we may construct pairwise disjoint closed balls $B_{n,1},\dots,B_{n,2^{n}}$ such that each $B_{n,k}$ has radius no larger than $1/n$ and contains two of the subsets $B_{n+1,1},\dots,B_{n+1,2^{n+1}}$. The set $\bigcap_{n\geq 1}\bigcup_{k=1}^{2^n}B_{n,k}$ is a Cantor set.
\end{proof}
\begin{theorem}
  Let $X$ be a metric space, and let $Y$ be a complete metric space. Let $A\subseteq X$ be any subset, and let $f\colon A\to Y$ be a continuous map. Then, there exists a $G_{\delta}$ $B\subseteq X$ such that $A\subseteq B\subseteq \overline{A}$ and a continuous extension $g\colon B\to Y$ of $f$.
\end{theorem}
\begin{proof}
  Since $ \overline{A} $ is $G_{\delta}$ in $X$, we may assume that $ \overline{A} = X $. For each $n\in \N$, let
  \begin{align*}
    B_n &= \bigcup\set{U\subseteq X : \text{$U$ is open and $\operatorname{diam}\left( f\left( A\cap U \right) \right)\leq 1/n$}}.
  \end{align*}
  Then, $B = \bigcap_{n\in \N}B_n$ is a $G_{\delta}$ subset of $X$ and $A\subseteq B$ by the continuity of $f$. For each $x\in B$, the intersection $\bigcap_{U} \overline{f\left( A\cap U \right)}$ over all the open neighborhoods $U$ of $x$ is a singleton, which we call $\set{y}$. Setting $g(x) = y$, this is our desired continuous extension.
\end{proof}
\begin{proposition}
  Let $X$ be a complete metrizable space, and let $Y$ be a nonempty subset of $X$. Then, $Y$ is a completely metrizable space if and only if $Y$ is a $G_{\delta}$ subset of $X$. The same is also true upon substituting ``Polish space'' for ``completely metrizable space.''
\end{proposition}
\begin{proof}
  We will prove the case for completely metrizable spaces. Let $d$ be a compatible metric on $X$. If $Y$ is open, define
  \begin{align*}
    d'\left( y,z \right) &= d\left( y,z \right) + \left\vert \frac{1}{\dist\left( y,X\setminus Y \right)} - \frac{1}{\dist\left( z,X\setminus Y \right)} \right\vert.
  \end{align*}
  If $Y$ is a $G_{\delta}$, then $Y = \bigcap_{n\geq 1}U_n$ for some open subsets $U_n$, so the diagonal embedding $Y\to \prod_{n\geq 1}U_n$ realizes $Y$ as a closed subspace of $\prod_{n\geq 1}U_n$, so $Y$ is a completely metrizable space.
\end{proof}
\begin{proposition}
  Every Polish space is homeomorphic to a $G_{\delta}$ subset of $\left[ 0,1 \right]^{\N}$.
\end{proposition}
\begin{proof}
  Let $X$ be a Polish space, $d$ a bounded compatible metric. Taking a dense subset $\set{p_n : n\in \N}$ in $X$, the map $X\mapsto \left[ 0,1 \right]^{\N}$ taking $x\mapsto \left( d\left( x,p_n \right) \right)_{n\in \N}$ is a homeomorphism of $X$ onto its image. So, $X$ is a $G_{\delta}$ subset of $\left[ 0,1 \right]^{\N}$.
\end{proof}
\begin{theorem}
  Let $\left( X,\mathcal{T} \right)$ be a Polish space. Let $A\in \mathcal{B}_X$. There exists a Polish topology $\mathcal{T}_A$ on $X$ that is finer than $\mathcal{T}$ with $\sigma\left( \mathcal{T}_A \right) = \mathcal{B}_X$ and $A$ is clopen in $\left( X,\mathcal{T}_A \right)$.
\end{theorem}
\begin{proof}
  Let $\mathcal{A}$ be the set of all $A\in \mathcal{B}_X$ for which the conclusion holds. We will show that $\mathcal{A}$ is a $\sigma$-algebra containing $\mathcal{T}_A$. Let $A\in \mathcal{T}$. Then, both $A$ and $X\setminus A$ are Polish spaces. If $\mathcal{T}_A$ is the disjoint union topology on $A\sqcup X\setminus A$, then $\mathcal{T}_A$ is Polish and $\mathcal{T}_A\supseteq \mathcal{T}$, with $\sigma\left( \mathcal{T}_A \right) = \mathcal{B}_X$. Moreover, $A$ is clopen in $\left( X,\mathcal{T}_A \right)$, so $\mathcal{T}\subseteq \mathcal{A}$. 

  Our goal now is to show that $\mathcal{A}$ is a $\sigma$-algebra. If $\set{A_n : n\in \N}$ is a sequence in $\mathcal{A}$, then the topology $\mathcal{T}_A$ generated by $\bigcup_{n\geq 1}\mathcal{T}_{A_n}$ is Polish and finer than $\mathcal{T}$. Since each $\mathcal{T}_{A_n}$ has a countable basis and is contained in $\mathcal{B}_X$, we have that $\mathcal{T}_A\subseteq \mathcal{B}_X$, so $\sigma\left( \mathcal{T}_A \right) = \mathcal{B}_X$. Moreover, $\bigcup_{n\geq 1}A_n$ is open in $\mathcal{T}_A$, so there is a Polish topology $\mathcal{T}_{A}'$ finer than $\mathcal{T}_A$ such that $\sigma\left( \mathcal{T}_A' \right) = \sigma\left( \mathcal{T}_A \right) = \mathcal{B}_X$, so $\mathcal{A}$ is closed under taking countable unions. Finally, since we may replace $A$ with $X\setminus A$ in the way we defined the topology $\mathcal{T}_A$, it follows that $\mathcal{A}$ is a $\sigma$-algebra.
\end{proof}
\begin{corollary}
  Let $\left( X,\mathcal{T} \right)$ be a Polish space, $Y$ a second countable topological space, and $f\colon X\to Y$ a Borel map. Then, there is a Polish topology $\mathcal{T}_f$ on $X$ such that $\sigma\left( \mathcal{T}_f \right) = \mathcal{B}_X$ and $f\colon \left( X,\mathcal{T}_f \right)\to Y$ is continuous.
\end{corollary}
\begin{theorem}[Schröder--Bernstein for Measurable Spaces]
  Let $\left( X,\mathcal{B} \right)$ and $\left( Y,\mathcal{C} \right)$ be measurable spaces. Let $f\colon X\to Y$ and $g\colon Y\to X$ be injections such that $f\left( \mathcal{B} \right)\subseteq \mathcal{C}$ and $g\left( \mathcal{C} \right)\subseteq \mathcal{B}$ such that each establishes an isomorphism of measurable spaces onto its image. Then, there exist $A\in \mathcal{B}$ and $B\in \mathcal{C}$ such that $f(A) = Y\setminus B$ and $g(B) = X\setminus A$.
\end{theorem}
\begin{proof}
  Let $B_1 = Y\setminus f(X)$ and $B = \bigcup_{n\geq 0}\left( f\circ g \right)^{n}\left(B_1\right)$. Set $A = X\setminus g(B)$. Then, $g(B) = X\setminus A$ and
  \begin{align*}
    f(A) &= f\left( X\setminus g(B) \right)\\
         &= f\left( X \right)\setminus \left( f\circ g \right)(B)\\
         &= Y\setminus \left( B_1\cup \left( f\circ g \right)(B) \right)\\
         &= Y\setminus B.
  \end{align*}
\end{proof}
\begin{lemma}
  Let $\set{\left( X_n,\mathcal{T}_n \right) : n\in \N}$ be a sequence of second countable topological spaces. Then, $ \sigma\left( \prod_{n=1}^{\infty}\mathcal{T}_n \right) = \prod_{n=1}^{\infty}\sigma\left( \mathcal{T}_n \right) $.
\end{lemma}
\begin{proof}
  We observe that $\pi_n\left( \sigma\left( \mathcal{T}_n \right) \right) = \sigma\left( \pi_n^{-1}\left( \mathcal{T}_n \right) \right)$ for each $n\in \N$, we have that
  \begin{align*}
    \prod_{n=1}^{\infty} \sigma\left( \mathcal{T}_n \right) &= \sigma\left( \bigcup_{n=1}^{\infty}\pi_n^{-1}\left( \mathcal{T}_n \right) \right).
  \end{align*}
  Letting $\mathcal{U}_n$ be a countable basis for $X_n$, we let $\mathcal{U}$ be the basis for $\prod_{n=1}^{\infty}\mathcal{T}_n$ given by cylinder subsets of the form $\prod_{k=1}^{n}\pi_k^{-1}\left( U_k \right)$. Then, we have
  \begin{align*}
    \sigma\left( \prod_{n=1}^{\infty}\mathcal{T}_n \right) &= \sigma\left( \mathcal{U} \right)\\
                                                           &= \sigma\left( \bigcup_{n=1}^{\infty}\pi_n^{-1}\left( \mathcal{T}_n \right) \right)\\
                                                           &= \prod_{n=1}^{\infty}\sigma\left( \mathcal{T}_n \right).
  \end{align*}
\end{proof}
\begin{lemma}
  The closed unit interval $\left[ 0,1 \right]$ with the Borel structure associated to the standard topology is isomorphic to the Cantor set with its standard Borel structure.
\end{lemma}
\begin{proof}
  Let $\varphi\colon \set{0,1}^{\N}\to [0,1]$ be given by
  \begin{align*}
    \varphi\left( \left( a_n \right)_n \right) &= \sum_{n=1}^{\infty} 2^{-n}a_n.
  \end{align*}
  We have that $\varphi$ is continuous and surjective. Let $Y$ be the set of points in $\left[ 0,1 \right]$ with more than one preimage under $\varphi$. Then, $Y$ and $\varphi^{-1}\left( Y \right)$ are countably infinite, so we may restrict $\varphi$ to $\set{0,1}^{\N}\setminus \varphi^{-1}\left( Y \right)$ to yield a homomorphism onto $\left[ 0,1 \right]\setminus Y$, hence a Borel isomorphism. Taking a bijection $\psi\colon \varphi^{-1}\left( Y \right)\to Y$ and stitching accordingly yields a Borel isomorphism $\set{0,1}^{\N}\to [0,1]$.
\end{proof}
\section{Factor Maps}%
\begin{definition}
  Two continuous actions $G\curvearrowright X$ and $G\curvearrowright Y$ by homeomorphisms are said to be topologically conjugate if there is an equivariant homeomorphism $\varphi\colon X\to Y$.

  Two probability measure-preserving actions $G\curvearrowright \left( X,\mu \right)$ and $G\curvearrowright \left( Y,\nu \right)$ on standard probability spaces are said to be measure conjugate if there are conull $G$-invariant subsets $X'\subseteq X$ and $Y'\subseteq Y$ and an equivariant measure isomorphism $\varphi\colon X'\to Y'$.
\end{definition}
\begin{definition}
  Let $G\curvearrowright X$ and $G\curvearrowright Y$ be continuous actions by homeomorphisms. We say the second is a factor of the first (and the first is an extension of the second) if there is an equivariant continuous surjection $\pi\colon X\to Y$.

  Let $G\curvearrowright \left( X,\mu \right)$ and $G\curvearrowright \left( Y,\nu \right)$ be probability measure-preserving actions. We say the second is a factor of the first if there is a $G$-invariant conull subset $X'\subseteq X$ and an equivariant map $\pi\colon X'\to Y$ such that $\pi_{\ast}\mu = \nu$. The extension is trivial if there is a $G$-invariant conull set $Y'\subseteq Y$ such that $\pi$ is a measure isomorphism between $X'$ and $Y'$.
\end{definition}
\nocite{kerr_and_li}
\printbibliography
\end{document}
